\begin{question}

  \begin{questionpart}
    Construct the transformation matrix that connects the $S_z$
    diagonal basis to the $S_x$ diagonal basis.
  \end{questionpart}

  \begin{questionpart}
    Show that your result is consistent with the general relation
    \begin{equation*}
      U = \sum_r \dyad{b^{(r)}}{a^{(r)}}
    \end{equation*}
  \end{questionpart}

\end{question}

\begin{purpose}
  He wants us to feel the connection between matrices and bras and
  kets.
\end{purpose}

\begin{answer}

  \begin{answerpart}{With LinAlg}
    So...the question is ambiguous.  I don't know if it wants $x
    \rightarrow z$ or $z \rightarrow x$.  So we'll just work it both
    ways. We'll start with a matrix that we can use in the $x$ basis,
    to get the $z$ basis. That is to say

    \begin{equation}
      \begin{split}
        \hat{O}_z\ket{\alpha_x} &= \ket{\alpha_z}\\
        \begin{bmatrix}
          O_{11} & O_{12}\\
          O_{21} & O_{22}\\
        \end{bmatrix}
        \begin{bmatrix}
          1 & 0\\
          0 & 1\\
        \end{bmatrix}
        &= 
        \begin{bmatrix}
          [x_{+,z}] & [x_{-,z}]\\
        \end{bmatrix}
      \end{split}
    \end{equation}

    where $\ket{\alpha_z}$ is $\ket{\alpha}$ expressed in the $z$
    basis and the same for the $x$ basis.  So, clearly all we have to
    do is write out the $\ket{x\pm}$ vectors in the $z$ basis and we
    win.

    \begin{equation}
      \hat{O}_z
      =
      \frac{1}{\sqrt{2}}
      \begin{bmatrix}
        1 & 1\\
        1 & -1\\
      \end{bmatrix}
    \end{equation}

    Drawing upon our experience in \ref{1.26:theQuestion}, we can do
    this one pretty easily too\footnote{Check out equation
    \ref{1.26:yzExpression} where we do this for y which differs only
    by a factor of $i$.}
    
    \begin{equation}
      \hat{O}_x
      =
      \frac{1}{\sqrt{2}}
      \begin{bmatrix}
        1 & 1\\
        1 & -1\\
      \end{bmatrix}
    \end{equation}

    So, it \emph{is} worth drawing our attention to the fact that the
    matrices are \emph{exactly} the same.  Let's take a look at why
    that is.

    \begin{equation}
      \begin{split}
        \hat{O}_z\ket{\alpha_x} &= \ket{\alpha_z}\\
        \hat{O}_x\ket{\alpha_z} &= \ket{\alpha_x}\\
        \hat{O}_z\left(\hat{O}_x\ket{\alpha_z}\right) &= \ket{\alpha_z}\\
        \hat{O}_z\hat{O}_x &= I\\
        \hat{O}_z &= \hat{O}_{x}^{-1}\\
      \end{split}
    \end{equation}

    Ladies and gentlement, it looks like we have ourselves a
    party... maybe not a hoppin' party but a party nonetheless.  Next
    up, let's take a look at the hermiticity of these things.  For the
    relations above to hold (which we'll show in a minute here) then
    the above relations also have to hold for $\bra{\alpha_{x/y}}$
    which means that

    \begin{equation}
      \begin{split}
        \bra{\alpha_x}\hat{O}_z^\dagger &= \bra{\alpha_z}\\
        \mel{\alpha_x}{\hat{O}_z^\dagger\hat{O}_z}{\alpha_x} &= \braket{\alpha_z}{\alpha_z}\\
        \mel{\alpha_x}{\hat{1}}{\alpha_x} &= \braket{\alpha_z}{\alpha_z}\\
        \hat{O}_z^\dagger\hat{O}_z &= I\\
        \hat{O}_x^\dagger\hat{O}_x &= I\\
        \hat{O}_x^\dagger &= \hat{O}_z\\
        \hat{O}_z^\dagger &= \hat{O}_x\\
      \end{split}
    \end{equation}

    It's worth noting that this relationship is a lot easier to show
    using bra/ket algebra with

    \begin{equation}
      \begin{split}
        \hat{O}_x &= \sum_i \dyad{x_i}{z_i}\\
        \hat{O}_z &= \sum_i \dyad{z_i}{x_i}\\
        &= \left(\sum_i \dyad{x_i}{z_i}\right)^\dagger\\
        &= \hat{O}_x^\dagger\\
      \end{split}
    \end{equation}

    In the general case, this is where the story ends.  In the
    \emph{specific} case, however, it's pretty easy to see that
    $\hat{O}_z^\dagger = \hat{O}_z = \hat{O}_x$.
    
  \end{answerpart}

  \begin{answerpart}{Bring in the Unitarians}
    Next up is showing compliance with the compliance requirement with
    bra/ket algebras.

    \begin{equation}
      \begin{split}
        O_z
        &= \sum_i \dyad{z_i}{x_i}\\
        &= \dyad{z_+}{x_+} + \dyad{z_-}{x_-}\\
        &\dot{=}
        \begin{bmatrix}
          \mel{z+}{O_z}{z+} & \mel{z+}{O_z}{z-}\\
          \mel{z-}{O_z}{z+} & \mel{z-}{O_z}{z-}\\
        \end{bmatrix}
        \\
        &=
        \begin{bmatrix}
          \mel{z+}{(\dyad{z_+}{x_+})}{z+} & \mel{z+}{(\dyad{z_+}{x_+})}{z-}\\
          \mel{z-}{(\dyad{z_-}{x_-})}{z+} & \mel{z-}{(\dyad{z_-}{x_-})}{z-}\\
        \end{bmatrix}
        \\
        &=
        \begin{bmatrix}
          \braket{x_+}{z+} & \braket{x_+}{z-}\\
          \braket{x_-}{z+} & \braket{x_-}{z-}\\
        \end{bmatrix}
        \\
        &=
        \frac{1}{\sqrt{2}}
        \begin{bmatrix}
          1 & 1\\
          1 & -1\\
        \end{bmatrix}
        \\
      \end{split}
    \end{equation}

    Because of the unitarian nature of the matrices we can either
    repeat that whole block for the other one, or just take it on
    faith that you can do that yourself...I believe in you...
    
  \end{answerpart}

\end{answer}
